\begin{pastexam}[label={px:ncu-112-dsa-04}]{中央資工所 112 學年度 複選題第 4 題}
（複選；每題 5 分，答錯 1 題倒扣 1 分）

Given the input list L = [28, 203, 16, 30, 123, 521, 63, 528, 210, 216, 941, 45]. Which of the
following statements is (are) correct?
\begin{enumerate}[label=(\Alph*)]
\item LSD radix sort is a non-comparative sorting algorithm.
\item At the end of the second pass of LSD Radix sort, the sixth element of the resulting chain is 123.
\item At the end of the third pass of LSD Radix sort, the sixth element of the resulting chain is 123.
\item At the end of the first pass of LSD Radix sort, the sixth element of the resulting chain is 63.
\end{enumerate}
\end{pastexam}

\begin{pastsoln}
答案是 (A)、(B)、(C)。

\textbf{(A) 對}。 LSD 基數排序只看每一位的數字把資料分到 $10$ 個桶子，從不拿兩筆資料互相比較大小，是
非比較式的排序，所以不受 $\Omega(n \log n)$ 比較下界的限制（資結書第~15.2.2~節）。

\textbf{追蹤}。 以 $10$ 為基數，三位數，第一趟看個位、第二趟看十位、第三趟看百位（不足三位的補 $0$）。
每一趟依原來的順序把數字接到對應桶子的尾端，再從桶子 $0$ 到 $9$ 串起來。
\begin{center}
\small
\begin{tabular}{lll}
\toprule
趟 & 非空的桶子 & 串起來的結果 \\
\midrule
1（個位） & 0: 30, 210；1: 521, 941；3: 203, 123, 63； & 30, 210, 521, 941, 203, \textbf{123}, \\
 & 5: 45；6: 16, 216；8: 28, 528 & 63, 45, 16, 216, 28, 528 \\
\midrule
2（十位） & 0: 203；1: 210, 16, 216；2: 521, 123, 28, 528； & 203, 210, 16, 216, 521, \textbf{123}, \\
 & 3: 30；4: 941, 45；6: 63 & 28, 528, 30, 941, 45, 63 \\
\midrule
3（百位） & 0: 16, 28, 30, 45, 63；1: 123；2: 203, 210, 216； & 16, 28, 30, 45, 63, \textbf{123}, \\
 & 5: 521, 528；9: 941 & 203, 210, 216, 521, 528, 941 \\
\bottomrule
\end{tabular}
\end{center}

\textbf{(B) 對}。 第二趟結束，第六個是 $123$。

\textbf{(C) 對}。 第三趟結束已經排好，比 $123$ 小的有 $16, 28, 30, 45, 63$ 五個，第六個是 $123$。

\textbf{(D) 錯}。 第一趟結束，第六個是 $123$；$63$ 是第七個（個位是 $3$ 的桶子裡，$63$ 排在 $203$、$123$
後面，因為輸入裡它比較晚出現）。
\end{pastsoln}
